\uuid{IJr6}
\titre{Logique}
\auteur{bodin ; croizat ; sacre}
\organisation{exo7}
\licence{CC BY-NC-SA 4.0}
\datecreate{2022-03-08}
\niveau{L1}
\module{Algèbre}
\chapitre{Logique, ensemble, raisonnement}
\sousChapitre{Logique}
\difficulte{2}
\theme{Logique}
\exo7qcmid{hilisit.logique.005}

\contenu{
\begin{question}
\qclassification{100.01}
Soit $\mathcal{P}$ une assertion vraie et $\mathcal{Q}$ une assertion fausse.
Quelles sont les assertions vraies ?
\begin{answers}
	\good{$\mathcal{P} \text{ ou }\mathcal{Q}$}
	\bad{$\mathcal{Q} \text{ et }\mathcal{P}$}
	\good{$\mathcal{P} \text{ ou }\text{non}(\mathcal{Q})$}
	\bad{$\mathcal{Q} \text{ ou }\text{non}(\mathcal{P})$}
\end{answers}
\begin{explanations}
"$\mathcal{P} \text{ ou }\mathcal{Q}$" est vraie car l'un des deux termes est vrai.\\
"$\mathcal{P} \text{ et }\mathcal{Q}$" est fausse car l'un des termes est faux.\\
"$\mathcal{P} \text{ ou }\text{non}(\mathcal{Q})$" est vraie car l'un des deux termes de part et d'autre du "ou" est vrai.\\
"$\mathcal{Q} \text{ ou }\text{non}(\mathcal{P})$" est fausse car les deux termes de part et d'autre du "ou" sont faux ("$\text{non}(\mathcal{P})$" est une assertion fausse puisque $\mathcal{P}$ est vraie).
\end{explanations}
\end{question}
}
